\subsection{A sharp source-mass order for four-state sources}
\label{sec:sharp-source-mass}
The lower construction makes rare source states increasingly scarce. We now
show that its square-root-logarithmic order is the worst possible order for
four inputs. No lower bound on conditional label probabilities is needed.
The constants in the universal bound are not claimed to be optimal.

\begin{lemma}[A scalar logarithmic moment inequality]
\label{lem:log-moment}
Let $M\ge1$ and $0\le z\le M$. With continuous conventions at zero,
\[
 z(\log z)^2\le(2+\log M)(z\log z-z+1).
\]
\end{lemma}
\begin{proof}
For $0<z\le1$, put $t=-\log z\ge0$. The stronger inequality with
coefficient two is equivalent to
$t^2\le2(e^t-1-t)$, which follows from the exponential series.
For $z\ge1$, put $t=\log z\ge0$. The inequality with coefficient $2+t$
is equivalent to
\[
 t-2+(t+2)e^{-t}\ge0.
\]
The left side is zero at zero and has derivative
$1-(t+1)e^{-t}\ge0$. Since $t\le\log M$, the result follows.
Continuity handles $z=0$.
\end{proof}

\begin{lemma}[Square-root entropy perturbation]
\label{lem:entropy-lipschitz}
Let $\mu=(\mu_1,\ldots,\mu_m)$ be a positive probability vector. For
$u\in[0,\infty)^m$, define
\[
 E_\mu u^2=\sum_i\mu_i u_i^2,\qquad
 \operatorname{Ent}_\mu(u^2)
 =\sum_i\mu_i u_i^2\log\frac{u_i^2}{E_\mu u^2},\qquad
 F_\mu(u)=\sqrt{\operatorname{Ent}_\mu(u^2)},
\]
where $F_\mu(0)=0$. Then, for all nonnegative $u,v$,
\begin{equation}
 |F_\mu(u)-F_\mu(v)|
 \le\sqrt{2+\log(1/\min_i\mu_i)}
       \left(\sum_i\mu_i(u_i-v_i)^2\right)^{1/2}.
 \label{eq:entropy-lipschitz}
\end{equation}
\end{lemma}
\begin{proof}
First take strictly positive $u$ with $F_\mu(u)>0$. Put
$q=E_\mu u^2$ and $z_i=u_i^2/q$. Then
$z_i\le1/\mu_i\le1/\min_j\mu_j$, and differentiation gives
\[
 \partial_iF_\mu(u)=\frac{\mu_i u_i\log z_i}{F_\mu(u)}.
\]
Writing $C_\mu=2+\log(1/\min_i\mu_i)$, the squared gradient norm dual
to the weighted Euclidean norm satisfies
\begin{align*}
 \sum_i\frac{(\partial_iF_\mu(u))^2}{\mu_i}
 &=\frac{q\sum_i\mu_i z_i(\log z_i)^2}
          {\operatorname{Ent}_\mu(u^2)}\\
 &\le\frac{q C_\mu\sum_i\mu_i(z_i\log z_i-z_i+1)}
          {\operatorname{Ent}_\mu(u^2)}=C_\mu,
\end{align*}
where Lemma~\ref{lem:log-moment} and $\sum_i\mu_i z_i=1$ give the
last two steps. Cauchy--Schwarz bounds the derivative along a line segment
by $\sqrt{C_\mu}$ times its weighted speed.

For strictly positive $u,v$, the segment joining them has zero entropy
exactly at constant vectors. It either lies entirely in the line of constant
vectors, where the claim is immediate, or intersects that line at at most
one point. Integrate the derivative bound on the remaining intervals and
use continuity at their endpoints. For general nonnegative $u,v$, apply the
result to $u+\epsilon\boldsymbol{1},v+\epsilon\boldsymbol{1}$ and send
$\epsilon\downarrow0$. Their difference is unchanged and $F_\mu$ is
continuous, including at zero.
\end{proof}

The lemma is an analytic tool; the result below uses its one-sided effect
on a regrouping operation. It avoids paying a logarithmic comparison
factor on the distortion already incurred by the old partition.

\begin{theorem}[Four-state source-mass bound]
\label{thm:mass-bound}
If $|\mathcal X|=4$ and $p_x\ge\gamma>0$, there is a deterministic
hierarchy satisfying, simultaneously for $k=2,3$,
\begin{equation}
 D(\cH_k)\le K_\gamma D_k^*,\qquad
 R(\cH_k)\le K_\gamma R_k^*,\qquad
 K_\gamma=2+2\sqrt{3+\log(1/\gamma)}.
 \label{eq:mass-bound}
\end{equation}
The inequalities include zero-optimum cases without division.
\end{theorem}
\begin{proof}
Put $A=D_2^*$, $r=D_3^*$, and $C_\gamma=2+\log(1/\gamma)$.
Let $\cP_2$ be an optimal two-cell partition, and let $\{i,j\}$ be a
pair of minimum merge cost $r$. If $i,j$ belong to the same cell of
$\cP_2$, the partition merging only this pair refines $\cP_2$, so both
flat optima are simultaneously attained.

Otherwise let $a=p_i\le p_j=b$, exchanging their names if necessary.
Let $S,T$ be the cells of $\cP_2$ containing $i,j$, respectively, and
write $A_S=D(S)$, $A_T=D(T)$; thus $A=A_S+A_T$.
We move $i$ into $T$. In the new cell $T\cup\{i\}$, normalized source
weights are at least $\gamma$, since its total mass is at most one.
For each label $y$, apply Lemma~\ref{lem:entropy-lipschitz} to the vector
$(\sqrt{P_x(y)})_{x\in T\cup\{i\}}$ and to the vector with its
$i$th coordinate replaced by $\sqrt{P_j(y)}$. Multiply each resulting
square-root inequality by the square root of the cell mass, then use the
Euclidean triangle inequality across label coordinates to obtain
\begin{equation}
 \sqrt{D(T\cup\{i\})}
 \le\sqrt{D_{\rm rep}}+
      \sqrt{C_\gamma a\,\He^2(P_i,P_j)},
 \label{eq:one-sided-perturbation}
\end{equation}
where $D_{\rm rep}$ is the distortion of the artificial cell in which
$P_i$ is replaced by $P_j$, and
$\He^2(P,Q)=\sum_y(\sqrt{P(y)}-\sqrt{Q(y)})^2$.
The entropy expression equals KL cluster distortion by expanding each
label coordinate, so no probabilistic encoding claim is made about the
artificial replacement.

To bound $D_{\rm rep}$, evaluate its KL objective at the original centroid
$P_T$. The posterior $P_j$ now has weight $a+b\le2b$; every other
weight is unchanged. Minimizing over centroids can only help. Therefore
$D_{\rm rep}\le2A_T$.
Also $\KL(P\Vert Q)\ge\He^2(P,Q)$: writing
$f(z)=z\log z-z+1$, the scalar difference
\[
 f(z)-(\sqrt z-1)^2
 =2\sqrt z\{\sqrt z\log\sqrt z-\sqrt z+1\}\ge0
\]
proves this after multiplication by $Q$ and summation. The mixture
centroid of a positive-weight pair contains its members' supports.
The unconstrained weighted Euclidean centroid then gives
\[
 r\ge\frac{ab}{a+b}\He^2(P_i,P_j),\qquad
 a\He^2(P_i,P_j)\le2r.
\]
Thus \eqref{eq:one-sided-perturbation} implies
\[
 D(T\cup\{i\})\le
    \left(\sqrt{2A_T}+\sqrt{2C_\gamma r}\right)^2.
\]
Removing $i$ from $S$ cannot increase its distortion, by the Jensen
aggregation identity. The total new distortion $t$ is consequently bounded by
\begin{align}
 t&\le A_S+
        \left(\sqrt{2A_T}+\sqrt{2C_\gamma r}\right)^2\nonumber\\
  &\le\left(\sqrt{2A}+\sqrt{2C_\gamma r}\right)^2
   \le4A+4C_\gamma r.
 \label{eq:regrouping-bound}
\end{align}
For the second inequality use $A_S+2A_T\le2A$ and $A_T\le A$.
The moved cells are unions of the three fine cells
$\{i,j\},\{\ell\},\{m\}$. If removing $i$ empties $S$, split the
remaining cell into two nonempty unions of these fine cells. Refinement
does not increase distortion. This constructs a two-cell partition
compatible with the minimum-pair three-cell partition.

If $r=0$, the first bound in \eqref{eq:regrouping-bound} gives
$t\le2A$, and the fine distortion is zero; this also covers $A=0$.
Suppose $r>0$ and put $u=A/r\ge1$. Retaining $\cP_2$ and refining
it once yields fine distortion at most $A$, hence a simultaneous factor
at most $u$. The minimum-pair hierarchy constructed above has factor at
most $4+4C_\gamma/u$. Choose the better hierarchy. The increasing function
$u$ and decreasing function $4+4C_\gamma/u$ meet at
$2+2\sqrt{1+C_\gamma}$, so
\[
 \min\{u,4+4C_\gamma/u\}\le2+2\sqrt{1+C_\gamma}=K_\gamma.
\]
This proves the excess-loss claim. Adding the common nonnegative baseline
$B_0$ preserves any multiplicative factor at least one, proving the
statement for total risk.
\end{proof}

\begin{corollary}[Sharp worst-case source-mass order]
\label{cor:sharp-mass-order}
For $0<\gamma\le1/4$, take the supremum of $\rho_{\rm det}$, or
of $\rho_{\rm stoch}$, over four-state sources with masses at least
$\gamma$, arbitrary finite target alphabets, and positive flat total-risk
benchmarks. Both suprema have order
$\Theta(\sqrt{\log(1/\gamma)})$ as $\gamma\downarrow0$.
\end{corollary}
\begin{proof}
Theorem~\ref{thm:mass-bound} gives the upper bound; every deterministic
hierarchy is an admissible stochastic refinement. For the lower bound,
set $L=\log(1/\gamma)$ in Theorem~\ref{thm:main}. Its rare masses equal
$\gamma$ and its anchor masses exceed $\gamma$ eventually. Its positive
full-support total risks give the required lower order in both settings.
\end{proof}

The order is sharp over all four-state sources; the leading constants
in Theorem~\ref{thm:main} remain sharp only for the specified family.
The universal bound does not identify the extremal distribution or the
optimal worst-case coefficient. Rare source states are necessary for an
unbounded four-state price; extreme label probabilities alone do not suffice.
